\documentclass[11pt, a4paper]{article}

\usepackage[T1]{fontenc}
\usepackage[utf8]{inputenc}
\usepackage{tgpagella}
\usepackage{microtype}
\usepackage[spanish, es-noshorthands]{babel}
\addto\captionsspanish{\renewcommand{\tablename}{Tabla}}

\usepackage{amsmath, amssymb, amsfonts}
\usepackage{graphicx}
\usepackage{booktabs}
\usepackage[most]{tcolorbox}
\usepackage[margin=2.5cm]{geometry}
\usepackage{fancyhdr}
\usepackage{tikz}
\usetikzlibrary{shapes.geometric, arrows.meta, positioning, calc, 3d}

\pagestyle{fancy}
\fancyhf{}
\setlength{\headheight}{16pt}
\lhead{\textbf{UCB Sede Tarija} | Ing. Mecatrónica}
\chead{Examen EC1}
\rhead{Robótica (IMT-342) W06S01}
\lfoot{Prof: M.Sc. Ing. Bernardo Quiroga Turdera}
\rfoot{Página \thepage}
\renewcommand{\headrulewidth}{0.4pt}
\renewcommand{\footrulewidth}{0.4pt}

\definecolor{ucbblue}{RGB}{0, 51, 102}

\begin{document}

\begin{center}
    \Large\textbf{\textcolor{ucbblue}{EC1 — Marcos, Transformaciones, Denavit–Hartenberg y Cinemática Directa}}\\[0.2cm]
\end{center}

\textbf{Nombre:} \rule{8cm}{0.15mm} \hfill \textbf{Fecha:} \rule{4cm}{0.15mm}

\vspace{0.4cm}
\begin{tcolorbox}[colback=gray!5, colframe=gray!50, title=\textbf{Condiciones de la Evaluación}]
    \begin{itemize}
        \item Examen escrito individual.
        \item No se permite el uso de calculadora ni formulario.
        \item Muestre su razonamiento y transformaciones intermedias; respuestas finales sin respaldo tendrán crédito limitado.
        \item Utilice la convención D-H \textbf{estándar} estudiada en clase.
        \item Si detecta una inconsistencia en su propio resultado, indíquelo y explique por qué.
        \item A menos que se especifique lo contrario, los ángulos están en radianes y las distancias en metros.
    \end{itemize}
\end{tcolorbox}

\section*{Problema 1 — Marcos y Transformaciones Homogéneas (20 pts)}
Se conocen las siguientes matrices de transformación:
\begin{equation*}
    {}^{0}T_1 = \begin{bmatrix} 0 & -1 & 0 & 2 \\ 1 & 0 & 0 & 1 \\ 0 & 0 & 1 & 0 \\ 0 & 0 & 0 & 1 \end{bmatrix}, \qquad
    {}^{1}T_2 = \begin{bmatrix} 1 & 0 & 0 & 1 \\ 0 & 1 & 0 & 0 \\ 0 & 0 & 1 & 0 \\ 0 & 0 & 0 & 1 \end{bmatrix}
\end{equation*}

\textbf{a) [4 pts]} Explique con palabras qué representa físicamente ${}^{0}T_1$.

\textbf{b) [6 pts]} Sin realizar toda la multiplicación primero, escriba la expresión correcta para la pose del marco 2 respecto al marco 0. Luego, calcule numéricamente ${}^{0}T_2$.

\textbf{c) [6 pts]} Calcule ${}^{1}T_0$.

\textbf{d) [4 pts]} Un estudiante propone que ${}^{0}T_2 = {}^1T_2{}^0T_1$. Explique por qué esta cadena no es correcta.

\section*{Problema 2 — Interpretación D-H y Transformación (25 pts)}
Una junta revoluta $i$ se describe con los parámetros D-H estándar:
\begin{equation*}
    \theta_i = q_i, \qquad d_i = L, \qquad a_i = 0, \qquad \alpha_i = -\frac{\pi}{2}
\end{equation*}

\textbf{a) [6 pts]} Para cada parámetro ($\theta_i, d_i, a_i, \alpha_i$), indique su significado geométrico en la convención D-H estándar.

\textbf{b) [4 pts]} Identifique cuál de los parámetros D-H es la variable articular y explique por qué.

\textbf{c) [11 pts]} Construya la matriz de transformación homogénea completa ${}^{i-1}T_i(q_i)$.

\textbf{d) [4 pts]} Evalúe su matriz resultante en la configuración $q_i = 0$.

\section*{Problema 3 — Cinemática Directa de un Manipulador 3R (45 pts)}
Un manipulador serial didáctico 3R se describe mediante la siguiente tabla D-H estándar:

\begin{center}
    \begin{tabular}{c c c c c}
    \toprule
    $i$ & $\theta_i$ & $d_i$ & $a_i$ & $\alpha_i$ \\ \midrule
    1 & $q_1$ & $1$ & $0$ & $+\pi/2$ \\
    2 & $q_2$ & $0$ & $2$ & $0$ \\
    3 & $q_3$ & $0$ & $1$ & $0$ \\ \bottomrule
    \end{tabular}
\end{center}

\begin{center}
\begin{tikzpicture}[scale=1.5, thick]
    % Visual E2 - 3R Chain Concept
    \draw[->, gray] (0,0) -- (1,0) node[right]{$x_0$};
    \draw[->, gray] (0,0) -- (0,1) node[above]{$z_0$};
    
    \coordinate (O) at (0,0);
    \coordinate (J1) at (0,1);
    \draw[ucbblue, line width=2pt] (O) -- (J1) node[midway, left]{$d_1=1$};
    \filldraw[black] (O) circle (1.5pt) node[below left]{$\{0\}$};
    \filldraw[black] (J1) circle (1.5pt) node[left]{$\{1\}$};
    
    % Following alpha = pi/2, the new Z1 points into the page (or outward), and X1 points right.
    % Link a2 moves along X1.
    \coordinate (J2) at (2,1);
    \draw[ucbblue!70, line width=2pt] (J1) -- (J2) node[midway, above]{$a_2=2$};
    \filldraw[black] (J2) circle (1.5pt) node[above]{$\{2\}$};
    
    % Link a3
    \coordinate (J3) at (3,1);
    \draw[ucbblue!40, line width=2pt] (J2) -- (J3) node[midway, above]{$a_3=1$};
    \filldraw[red] (J3) circle (1.5pt) node[right]{$\{3\}$};
\end{tikzpicture}
\end{center}

Para la configuración $q_1 = 0, \quad q_2 = \pi/2, \quad q_3 = -\pi/2$, determine la pose del marco 3 respecto al marco 0.

\textbf{a) [15 pts]} Construya ${}^{0}T_1$, ${}^{1}T_2$ y ${}^{2}T_3$.

\textbf{b) [5 pts]} Escriba la cadena de transformación correcta para ${}^{0}T_3$.

\textbf{c) [13 pts]} Calcule numéricamente ${}^{0}T_3$.

\textbf{d) [6 pts]} Extraiga ${}^{0}p_3$ y ${}^{0}R_3$. Explique qué representa físicamente cada uno.

\textbf{e) [6 pts]} Realice una verificación matemática o física significativa de su resultado.

\section*{Problema 4 — Verificación y Diagnóstico de Errores (10 pts)}
Un robot planar 2R tiene longitudes de eslabón $a_1 = 2$ y $a_2 = 1$. Un estudiante reporta que la posición del efector final es:
\begin{equation*}
    p = \begin{bmatrix} 4 \\ 0 \\ 0 \end{bmatrix}
\end{equation*}

\textbf{a) [5 pts]} Sin resolver todo el problema de cinemática directa nuevamente, determine si este resultado puede ser físicamente correcto. Justifique su respuesta.

\textbf{b) [5 pts]} Otro estudiante obtiene la siguiente matriz de orientación:
\begin{equation*}
    R = \begin{bmatrix} 1 & 0 & 0 \\ 0 & 2 & 0 \\ 0 & 0 & 1 \end{bmatrix}
\end{equation*}
Dé una razón matemática por la cual esta no puede ser una matriz de rotación válida.

\end{document}
